\section{Formal Model}
\label{sec:formal}

This section extends the formal framework of BTV~\cite{btv2026} to the
persistence layer. We introduce two new definitions, one axiom governing the
trust model, and the central theorem of this paper.

\subsection{Inherited Definitions}

\begin{definition}[Materialized Verdict~{\normalfont\cite{btv2026}, Def.~4.1}]
\label{def:materialized}
A verdict $V$ is \emph{materialized} if and only if there exist an evidence
token~$E$, a compliance token~$C$, a decision outcome $d \in \{\Allow,
\Deny\}$, and an explanation string~$x$ such that
$V = \mathtt{Verdict{::}new}(E, C, d, x)$.
The function $\mathtt{Verdict{::}new}$ is the sole constructor of type
$\mathtt{Verdict}$; no other term of that type is well-typed in Safe Rust.
\end{definition}

\subsection{New Definitions}

\begin{definition}[Delivered Decision]
\label{def:delivered}
A materialized verdict $V$ is \emph{delivered} if the program has invoked
$\mathtt{DeliveryToken{::}deliver}$, transmitting the decision payload across
the API boundary to the requesting party.
\end{definition}

\begin{definition}[Ephemeral Verdict]
\label{def:ephemeral}
A verdict $V$ is \emph{ephemeral} if and only if $V$ is materialized
(Definition~\ref{def:materialized}), delivered (Definition~\ref{def:delivered}),
but no valid $\mathtt{InclusionReceipt}$~$R$ exists at the time of delivery
such that $R$ cryptographically attests to the inclusion of $V$ in a public
Transparency Log. Formally, let $\mathcal{E}$ denote the set of ephemeral
verdicts; we prove $\mathcal{E} = \emptyset$ by construction
(Theorem~\ref{thm:persistence-enclosure}).
\end{definition}

\subsection{Threat Model}

We model a \emph{System Operator} adversary with the following capabilities:
(i)~full control over the application process memory and deployment
infrastructure; (ii)~ability to modify, delay, or drop arbitrary network
requests; (iii)~no access to the private signing key $K_{\mathit{priv}}$ of the
Transparency Log authority.

\begin{axiom}[Log Authority]
\label{ax:log-authority}
Let $K_{\mathit{priv}}$ be the Ed25519 private signing key of the Transparency
Log. We assume $K_{\mathit{priv}}$ is not accessible to the System Operator.
Consequently, the operator cannot produce a valid
$\mathtt{InclusionReceipt}$ without a genuine submission to the Log.
\end{axiom}

\paragraph{Scope of the guarantee.}
Axiom~\ref{ax:log-authority} holds under the assumption of an honest Log key.
If $K_{\mathit{priv}}$ is compromised, a malicious operator could forge
receipts. This boundary is analogous to the \texttt{unsafe} boundary in
BTV~\cite{btv2026}: the type system enforces structural correctness; the
security of cryptographic infrastructure is a separate trust assumption, just
as Safe Rust relies on the absence of \texttt{unsafe} blocks. Detection of log
misbehaviour relies on independent monitors (the Gossip model of CT~\cite{rfc6962}),
ensuring that suppression remains detectable even if the operator controls
the log server infrastructure.

\subsection{The Persistence Enclosure Theorem}

\begin{theorem}[Persistence Enclosure]
\label{thm:persistence-enclosure}
In any Safe Rust program $P$ whose crate graph includes \texttt{btv-transparency},
the set of Ephemeral Verdicts $\mathcal{E}(P) = \emptyset$.
Equivalently, $\mathcal{E}(P) \neq \emptyset \Rightarrow P$ does not type-check.
\end{theorem}

\begin{proof}[Proof sketch]
By the type law~(\ref{eq:delivery-law}), the sole function that crosses the
API boundary is $\mathtt{DeliveryToken{::}deliver}$, which consumes a
$\mathtt{DeliveryToken}$ by value. The sole constructor of
$\mathtt{DeliveryToken}$ is $\mathtt{seal}(V, R)$, which requires both a
$\mathtt{Verdict}$~$V$ and an $\mathtt{InclusionReceipt}$~$R$.
By Axiom~\ref{ax:log-authority}, a valid $R$ can only be produced by the Log
Authority in response to a genuine submission. The encapsulation perimeter of
$\mathtt{InclusionReceipt}$ (private fields, $\mathtt{pub(crate)}$ constructor,
$\mathtt{\#[must\_use]}$) mirrors the three-case proof of the Constitutional
Enclosure Theorem~\cite{btv2026}: struct-literal forging (E0451), external
consumption (E0624), and silent drop (unused-\texttt{must\_use}) are all
ill-typed in Safe Rust. Therefore, no delivery path exists without a valid~$R$,
and $\mathcal{E}(P) = \emptyset$. The compile-fail test suite
(Section~\ref{sec:impl}) verifies all attack cases exhaustively.
\end{proof}

The full case-by-case proof, mirroring the structure of~\cite{btv2026}
Section~4.4, is given in Section~\ref{sec:impl}.
